\documentclass[10pt,a4paper]{amsart}
\usepackage{amsmath,amssymb,amsthm,graphicx}
\usepackage{verbatim}
\newenvironment{reference}{\comment}{\endcomment}
%\newenvironment{reference}{}{}
\newenvironment{slogan}{\comment}{\endcomment}
\newenvironment{history}{\comment}{\endcomment}

% For commutative diagrams we use Xy-pic
\usepackage[all]{xy}

% We use 2cell for 2-commutative diagrams.
\xyoption{2cell}
% Theorem environments.
%
\theoremstyle{plain}
\newtheorem{theorem}[subsection]{Theorem}
\newtheorem{proposition}[subsection]{Proposition}
\newtheorem{lemma}[subsection]{Lemma}

\theoremstyle{definition}
\newtheorem{definition}[subsection]{Definition}
\newtheorem{example}[subsection]{Example}
\newtheorem{exercise}[subsection]{Exercise}
\newtheorem{situation}[subsection]{Situation}

\theoremstyle{remark}
\newtheorem{remark}[subsection]{Remark}
\newtheorem{remarks}[subsection]{Remarks}

\numberwithin{equation}{subsection}

% Macros
%
\def\lim{\mathop{\mathrm{lim}}\nolimits}
\def\colim{\mathop{\mathrm{colim}}\nolimits}
\def\Spec{\mathop{\mathrm{Spec}}}
\def\Hom{\mathop{\mathrm{Hom}}\nolimits}
\def\Ext{\mathop{\mathrm{Ext}}\nolimits}
\def\SheafHom{\mathop{\mathcal{H}\!\mathit{om}}\nolimits}
\def\SheafExt{\mathop{\mathcal{E}\!\mathit{xt}}\nolimits}
\def\Sch{\mathit{Sch}}
\def\Mor{\mathop{\mathrm{Mor}}\nolimits}
\def\Ob{\mathop{\mathrm{Ob}}\nolimits}
\def\Sh{\mathop{\mathit{Sh}}\nolimits}
\def\NL{\mathop{N\!L}\nolimits}
\def\CH{\mathop{\mathrm{CH}}\nolimits}
\def\proetale{{pro\text{-}\acute{e}tale}}
\def\etale{{\acute{e}tale}}
\def\QCoh{\mathit{QCoh}}
\def\Ker{\mathop{\mathrm{Ker}}}
\def\Im{\mathop{\mathrm{Im}}}
\def\Coker{\mathop{\mathrm{Coker}}}
\def\Coim{\mathop{\mathrm{Coim}}}

% Boxtimes
%
\DeclareMathSymbol{\boxtimes}{\mathbin}{AMSa}{"02}

%
% Macros for moduli stacks/spaces
%
\def\QCohstack{\mathcal{QC}\!\mathit{oh}}
\def\Cohstack{\mathcal{C}\!\mathit{oh}}
\def\Spacesstack{\mathcal{S}\!\mathit{paces}}
\def\Quotfunctor{\mathrm{Quot}}
\def\Hilbfunctor{\mathrm{Hilb}}
\def\Curvesstack{\mathcal{C}\!\mathit{urves}}
\def\Polarizedstack{\mathcal{P}\!\mathit{olarized}}
\def\Complexesstack{\mathcal{C}\!\mathit{omplexes}}
% \Pic is the operator that assigns to X its picard group, usage \Pic(X)
% \Picardstack_{X/B} denotes the Picard stack of X over B
% \Picardfunctor_{X/B} denotes the Picard functor of X over B
\def\Pic{\mathop{\mathrm{Pic}}\nolimits}
\def\Picardstack{\mathcal{P}\!\mathit{ic}}
\def\Picardfunctor{\mathrm{Pic}}
\def\Deformationcategory{\mathcal{D}\!\mathit{ef}}
\usepackage{cleveref}
\crefname{originallemma}{original lemma}{original lemmas}
\Crefname{originallemma}{Original lemma}{Original lemmas}
\makeatletter
\newlabel{algebra-lemma-colimit-category-fp-algebras}{{05N9}{source}{}{sourcepointer-1}{}}
\newlabel{algebra-lemma-colimit-category-fp-algebras@cref}{{[originallemma][0][]05N9}{[1][0][]source}{}{}{}}
\makeatother
\crefname{originallemma}{original lemma}{original lemmas}
\Crefname{originallemma}{Original lemma}{Original lemmas}
\makeatletter
\newlabel{limits-lemma-descend-quasi-finite}{{094M}{source}{}{sourcepointer-2}{}}
\newlabel{limits-lemma-descend-quasi-finite@cref}{{[originallemma][0][]094M}{[1][0][]source}{}{}{}}
\makeatother
\crefname{originallemma}{original lemma}{original lemmas}
\Crefname{originallemma}{Original lemma}{Original lemmas}
\makeatletter
\newlabel{lemma-generically-finite}{{03S1}{source}{}{sourcepointer-3}{}}
\newlabel{lemma-generically-finite@cref}{{[originallemma][0][]03S1}{[1][0][]source}{}{}{}}
\makeatother
\crefname{originallemma}{original lemma}{original lemmas}
\Crefname{originallemma}{Original lemma}{Original lemmas}
\makeatletter
\newlabel{morphisms-lemma-finite-flat}{{02KB}{source}{}{sourcepointer-4}{}}
\newlabel{morphisms-lemma-finite-flat@cref}{{[originallemma][0][]02KB}{[1][0][]source}{}{}{}}
\makeatother
\crefname{originallemma}{original lemma}{original lemmas}
\Crefname{originallemma}{Original lemma}{Original lemmas}
\makeatletter
\newlabel{morphisms-lemma-composition-finite-locally-free}{{02KC}{source}{}{sourcepointer-5}{}}
\newlabel{morphisms-lemma-composition-finite-locally-free@cref}{{[originallemma][0][]02KC}{[1][0][]source}{}{}{}}
\makeatother
\crefname{originallemma}{original lemma}{original lemmas}
\Crefname{originallemma}{Original lemma}{Original lemmas}
\makeatletter
\newlabel{morphisms-lemma-massage-finite}{{03HW}{source}{}{sourcepointer-6}{}}
\newlabel{morphisms-lemma-massage-finite@cref}{{[originallemma][0][]03HW}{[1][0][]source}{}{}{}}
\makeatother
\title{A quasi-finite finitely presented morphism of affine schemes is decomposed by a finite locally free cover}
\author{The Stacks Project Authors}
\date{stacks-project:etale-cohomology.tex:local-snapshot-20261011; GNU FDL 1.2 or later}
\begin{document}
\maketitle
\noindent\textit{Editorial note.} All delivered mathematical text is the author's own, copied from the pinned Stacks Project chapter file etale-cohomology.tex; no statement or proof step was rewritten, reordered or removed.\par\medskip
\noindent\textit{Editorial note.} The proof is the author's complete argument; the parenthetical '(Details omitted.)' at source line 12114 concerns the standard descent of an algebra along a filtered colimit, whose inputs are named, and the remaining reductions are written out.\par\medskip
\noindent\textit{Editorial note.} Every labelled reference occurring inside the counted spans is carried in assembly.external\_references as a non-load-bearing navigation pointer (same-chapter ones at their definition site, cross-chapter ones at the invocation inside the counted span), so no load-bearing mathematics is replaced by a stub. No bibliography citation occurs inside the counted spans.\par\medskip
\noindent\textit{Editorial note.} The checked chapter list for same-name collisions is the set of stacks chapters staged in this repository: chow, cohomology, divisors, etale-cohomology, flat, more-algebra, more-morphisms, morphisms, spaces-more-morphisms, varieties together with this card's chapter etale-cohomology.\par\medskip
\noindent\textit{Editorial note.} The preamble is the chapter's own native preamble excerpt already used for this source: the Xy-pic setup (native preamble lines 14-18) and the native theorem declarations with the mathematical macros (native preamble lines 154-217).\par\medskip
\noindent\textit{Allowed standard prerequisites (editorial).}\begin{itemize}
\item Descent of quasi-finiteness. Limits, Lemma limits-lemma-descend-quasi-finite, together with Algebra, Lemma algebra-lemma-colimit-category-fp-algebras. Reduces the statement to finite type algebras over the integers, that is, to the Noetherian case.
\item Generic finiteness and Noetherian induction. Same chapter, Lemma lemma-generically-finite. Produces the partition into strata on which the morphism becomes finite.
\item Structure of a finite morphism and base change of finite locally free morphisms. Morphisms, Lemmas morphisms-lemma-massage-finite, morphisms-lemma-finite-flat and morphisms-lemma-composition-finite-locally-free. Supplies the model in which the finite morphism is a finite union of closed subschemes mapping isomorphically, and keeps finite local freeness under composition.
\end{itemize}\medskip
\ifdefined\hypertarget\hypertarget{sourcepointer-1}{}\fi
\noindent\textit{Source reference.} Original lemma 05N9: sources/stacks-etale-cohomology/pinned-original-source.tex, source lines 12110--12110. Navigation pointer invoked at proof line 12110; the counted span contains the pointer only and the cited result is not reproduced inside it, so the pointer is registered as a non-load-bearing navigation stub in the established form for Stacks material. No bibliography citation is involved.\par\medskip
\ifdefined\hypertarget\hypertarget{sourcepointer-2}{}\fi
\noindent\textit{Source reference.} Original lemma 094M: sources/stacks-etale-cohomology/pinned-original-source.tex, source lines 12112--12112. Navigation pointer invoked at proof line 12112; the counted span contains the pointer only and the cited result is not reproduced inside it, so the pointer is registered as a non-load-bearing navigation stub in the established form for Stacks material. No bibliography citation is involved.\par\medskip
\ifdefined\hypertarget\hypertarget{sourcepointer-3}{}\fi
\noindent\textit{Source reference.} Original lemma 03S1: sources/stacks-etale-cohomology/pinned-original-source.tex, source lines 12069--12069. Navigation pointer invoked at proof line 12118; the counted span contains the pointer only and the cited result is not reproduced inside it, so the pointer is registered as a non-load-bearing navigation stub in the established form for Stacks material. No bibliography citation is involved.\par\medskip
\ifdefined\hypertarget\hypertarget{sourcepointer-4}{}\fi
\noindent\textit{Source reference.} Original lemma 02KB: sources/stacks-etale-cohomology/pinned-original-source.tex, source lines 12136--12136. Navigation pointer invoked at proof line 12136; the counted span contains the pointer only and the cited result is not reproduced inside it, so the pointer is registered as a non-load-bearing navigation stub in the established form for Stacks material. No bibliography citation is involved.\par\medskip
\ifdefined\hypertarget\hypertarget{sourcepointer-5}{}\fi
\noindent\textit{Source reference.} Original lemma 02KC: sources/stacks-etale-cohomology/pinned-original-source.tex, source lines 12139--12139. Navigation pointer invoked at proof line 12139; the counted span contains the pointer only and the cited result is not reproduced inside it, so the pointer is registered as a non-load-bearing navigation stub in the established form for Stacks material. No bibliography citation is involved.\par\medskip
\ifdefined\hypertarget\hypertarget{sourcepointer-6}{}\fi
\noindent\textit{Source reference.} Original lemma 03HW: sources/stacks-etale-cohomology/pinned-original-source.tex, source lines 12146--12146. Navigation pointer invoked at proof line 12146; the counted span contains the pointer only and the cited result is not reproduced inside it, so the pointer is registered as a non-load-bearing navigation stub in the established form for Stacks material. No bibliography citation is involved.\par\medskip
\begin{lemma}
\label{lemma-decompose-quasi-finite-morphism}
Let $f : Y \to X$ be a quasi-finite and finitely presented
morphism of affine schemes.
\begin{enumerate}
\item There exists a surjective morphism of affine schemes $X' \to X$ and a
closed subscheme $Z' \subset Y' = X' \times_X Y$ such that
\begin{enumerate}
\item $Z' \subset Y'$ is a thickening, and
\item $Z' \to X'$ is a finite \'etale morphism.
\end{enumerate}
\item There exists a finite partition $X = \coprod X_i$ by
locally closed, constructible, affine strata, and surjective finite locally
free morphisms $X'_i \to X_i$ such that the reduction of
$Y'_i = X'_i \times_X Y \to X'_i$ is isomorphic to
$\coprod_{j = 1}^{n_i} (X'_i)_{red} \to (X'_i)_{red}$ for some $n_i$.
\end{enumerate}
\end{lemma}

\begin{proof}
Setting $X' = \coprod X'_i$ we see that (2) implies (1).
Write $X = \Spec(A)$ and $Y = \Spec(B)$. Write $A$ as a filtered colimit
of finite type $\mathbf{Z}$-algebras $A_i$. Since $B$ is an $A$-algebra of
finite presentation, we see that there exists $0 \in I$ and a
finite type ring map $A_0 \to B_0$ such that $B = \colim B_i$ with
$B_i = A_i \otimes_{A_0} B_0$, see
Algebra, Lemma \ref{algebra-lemma-colimit-category-fp-algebras}.
For $i$ sufficiently large we see that $A_i \to B_i$ is
quasi-finite, see Limits, Lemma \ref{limits-lemma-descend-quasi-finite}.
Thus we reduce to the case of finite type algebras over $\mathbf{Z}$,
in particular we reduce to the Noetherian case. (Details omitted.)

\medskip\noindent
Assume $X$ and $Y$ Noetherian. In this case any locally closed
subset of $X$ is constructible. By Lemma \ref{lemma-generically-finite}
and Noetherian induction we see that
there is a finite partition $X = \coprod X_i$ of $X$
by locally closed strata such that $Y \times_X X_i \to X_i$ is finite.
We can refine this partition to get affine strata.
Thus after replacing $X$ by $X' = \coprod X_i$ we may assume
$Y \to X$ is finite.

\medskip\noindent
Assume $X$ and $Y$ Noetherian and $Y \to X$ finite.
Suppose that we can prove (2) after base change by a surjective,
flat, quasi-finite morphism $U \to X$. Thus we have a partition
$U = \coprod U_i$ and finite locally free morphisms $U'_i \to U_i$
such that $U'_i \times_X Y \to U'_i$ is isomorphic to
$\coprod_{j = 1}^{n_i} (U'_i)_{red} \to (U'_i)_{red}$ for some $n_i$.
Then, by the argument in the previous paragraph, we can find a
partition $X = \coprod X_j$ with locally closed affine strata such that
$X_j \times_X U_i \to X_j$ is finite for all $i, j$. By
Morphisms, Lemma \ref{morphisms-lemma-finite-flat}
each $X_j \times_X U_i \to X_j$ is finite locally free.
Hence $X_j \times_X U'_i \to X_j$ is finite locally free
(Morphisms, Lemma \ref{morphisms-lemma-composition-finite-locally-free}).
It follows that $X = \coprod X_j$ and $X_j' = \coprod_i X_j \times_X U'_i$
is a solution for $Y \to X$. Thus it suffices to prove
the result (in the Noetherian case) after a surjective flat quasi-finite
base change.

\medskip\noindent
Applying Morphisms, Lemma \ref{morphisms-lemma-massage-finite}
we see we may assume that $Y$ is a closed subscheme of an
affine scheme $Z$ which is (set theoretically) a finite union
$Z = \bigcup_{i \in I} Z_i$ of closed subschemes mapping isomorphically
to $X$. In this case we will find a finite partition of $X = \coprod X_j$
with affine locally closed strata that works (in other words $X'_j = X_j$).
Set $T_i = Y \cap Z_i$. This is a closed subscheme of $X$.
As $X$ is Noetherian we can find a finite partition of $X = \coprod X_j$
by affine locally closed subschemes, such that each
$X_j \times_X T_i$ is (set theoretically) a union of strata $X_j \times_X Z_i$.
Replacing $X$ by $X_j$ we see that we may assume $I = I_1 \amalg I_2$
with $Z_i \subset Y$ for $i \in I_1$ and $Z_i \cap Y = \emptyset$ for
$i \in I_2$. Replacing $Z$ by $\bigcup_{i \in I_1} Z_i$ we see that we
may assume $Y = Z$.
Finally, we can replace $X$ again by the members of a partition
as above such that for every $i, i' \subset I$ the intersection
$Z_i \cap Z_{i'}$ is either empty or (set theoretically) equal
to $Z_i$ and $Z_{i'}$. This clearly means that $Y$ is (set theoretically)
equal to a disjoint union of the $Z_i$ which is what we wanted to show.
\end{proof}

\end{document}
